\chapter{Kurzfassung}

Large Language Models (LLMs) übernehmen zunehmend eine vermittelnde Rolle bei der Produktsuche, und die von ihnen zitierten Quellen beeinflussen unmittelbar die Sichtbarkeit von Marken und Publishern. Aufgrund der Black-Box-Natur von LLMs fehlt der aufkommenden Praxis der Generative Engine Optimization (GEO) das fundierte Verständnis, das die Suchmaschinenoptimierung (SEO) über Jahrzehnte aufgebaut hat. Diese Arbeit untersucht, wie ChatGPT und Gemini Webquellen auswählen und zitieren, wenn sie Produktempfehlungsanfragen beantworten. Wir entwickelten eine Methodik, um einen gegebenen Prompt an beide LLMs zu senden, die Antwort zu erfassen und zu extrahieren, welche Suchanfragen die Modelle im Hintergrund stellen und welche strukturellen Merkmale die zitierten Seiten aufweisen. Wir führten 79~Prompts jeweils dreimal über zwei ChatGPT-Abonnementstufen (Plus und Business) sowie Gemini aus, was 711~Durchläufe ergab. Anschließend untersuchten wir, wie häufig LLM-Zitationen in den Bing- und Google-Suchergebnissen für dieselben Anfragen erscheinen und auf welcher Position sie ranken. Um die Treiber der Quellenauswahl zu verstehen, profilierten wir die strukturellen Eigenschaften von circa 6\,500~URLs und verglichen, was in den Suchergebnissen verfügbar war mit dem, was die LLMs tatsächlich zitierten. Zusätzlich überprüften wir, ob die Modelle Produktdetails korrekt wiedergeben, indem wir über 1\,100~Einzelaussagen aus zitierten Quellen auditierten.

77,7\,\% der von Gemini zitierten Links fanden sich in den Google-Suchergebnissen für dieselben Anfragen. Bei ChatGPT Business erschienen 81,3\,\% in den Bing-Ergebnissen, jedoch nur 27,8\,\% bei Google, während die Zitationen von Plus in beiden Suchmaschinen auftauchten (67,6\,\% Bing, 64,6\,\% Google), was auf eine Nutzung mehrerer Suchanbieter hindeutet. Über alle Systeme hinweg erschienen 16--22\,\% der Zitationen in keinem Suchergebnis für die jeweiligen Anfragen. Diese stammten überwiegend von Wikipedia, Reddit, App-Stores und großen Technikpublikationen. Sowohl ChatGPT als auch Gemini bevorzugten bei der Auswahl aus Listicles (Listenartikel wie "`Top~10 Best\ldots"') Seiten mit Tabellen, nummerierten Listen und Aufzählungszeichen und zeigten eine Präferenz für kürzlich veröffentlichte Inhalte. Produkte, die am Anfang eines Listicles platziert waren, wurden deutlich häufiger zitiert als solche weiter unten. Wenn die Modelle eine Quelle zitierten, gaben sie deren Inhalt in der Regel korrekt wieder.

Unsere Ergebnisse zeigen, dass ein gutes Suchmaschinenranking eine Voraussetzung für LLM-Sichtbarkeit ist, und verdeutlichen zugleich, wo GEO von traditioneller SEO abweicht.
